% Compile twice: lualatex 1500023.tex
\documentclass[11pt]{article}
\usepackage{fontspec}
\setmainfont{Latin Modern Roman}
\usepackage{amsmath,amssymb,mathtools,unicode-math}
\setmathfont{Latin Modern Math}
\usepackage{xurl,hyperref}
\hypersetup{colorlinks=true,urlcolor=blue,pdftitle={UnsolvedMath Problem 1500023}}
\newcommand{\sref}[2]{\hyperlink{source:#1:#2}{[S#2]}}
\newcommand{\cataloguescope}[1]{\paragraph{Mathematical question.}#1}
\begin{document}
% UnsolvedMath ID 1500023; source catalogue code AMR-014-0023
\setcounter{section}{1500022}
\section{Algebraic Stories — Exterior algebras}\label{1500023}
{\small\sffamily UnsolvedMath ID 1500023 · Algebra}
\subsection{Definitions and mathematical statement}

\paragraph{Shared notation.}$\mathbb N=\{1,2,\ldots\}$, and $\mathbb Z,\mathbb Q,\mathbb R,\mathbb C$ denote the integers, rationals, reals and complex numbers. Any other conventions are specified locally.


Let $E=\bigwedge_{\mathbb C}\mathbb C^n$ be the exterior algebra on $n$ degree-one generators, with graded pieces $E_i=\bigwedge^i\mathbb C^n$. For a homogeneous element $f\in E_d$ define
\[\operatorname{Ann}_E(f)=\{a\in E:af=0\},\qquad (f)=Ef.
\]
If $d$ is odd then $f^2=0$, hence $(f)\subseteq\operatorname{Ann}_E(f)$.
A property holds for a generic form in $E_d$ if it holds on a nonempty Zariski-open subset of $E_d$.
\paragraph{Question in the published source.}
For every $n\geq1$ and every odd integer $d$ with $1\leq d\leq n$, does a generic form $f\in E_d$ satisfy
\[(\operatorname{Ann}_E(f))_i=(f)_i\qquad\text{for every integer }0\leq i<(n-d)/2?\] \sref{1500023}{2}
\subsection{Short English statement}
Does a generic odd exterior form have only the annihilators generated by itself below the stated degree threshold?
\subsection{Source scope and status}
The published Problem Q explicitly assumes a generic form. The earlier arXiv v1, Problem P, omits that hypothesis; the previous catalogue wording asked about every form. These formulations are different. \sref{1500023}{2}\sref{1500023}{3}

The every-form variant is false: for $n=6$, $d=3$, and $f=e_1\wedge e_2\wedge e_3$, one has $e_1f=0$ but $e_1\notin(f)_1=0$, with $1<(6-3)/2$. This nongeneric example does not disprove the question above. The reviewed report also gives positive generic cases, including $n\leq d+4$ and cubic forms with $n\leq9$, consistent with the cited literature. It does not resolve the generic question for all parameters. \sref{1500023}{4}\sref{1500023}{5}
\subsection{Sources}
\begin{description}
\item[\hypertarget{source:1500023:1}{S1}] ULAM.AI and the UnsolvedMath Contributors, UnsolvedMath: A Curated Collection of Open Mathematics Problems, record 1500023 (AMR-014-0023). CC BY 4.0. \url{https://huggingface.co/datasets/ulamai/UnsolvedMath}.
\item[\hypertarget{source:1500023:2}{S2}] R. Fröberg, S. Lundqvist, A. Oneto and B. Shapiro, \emph{Algebraic Stories from One and from the Other Pockets}, Arnold Mathematical Journal 4 (2018), 137--160, §4.4, Problem Q, p.~155. \url{https://doi.org/10.1007/s40598-018-0088-z}.
\item[\hypertarget{source:1500023:3}{S3}] The same authors, earlier arXiv version, §4.4, Problem P. \url{https://arxiv.org/abs/1801.01692v1}.
\item[\hypertarget{source:1500023:4}{S4}] S. Lundqvist and L. Nicklasson, \emph{On generic principal ideals in the Exterior algebra}, arXiv:1803.03563v1 (2018), Propositions 4.2, 4.4 and 4.5. \url{https://arxiv.org/abs/1803.03563v1}.
\item[\hypertarget{source:1500023:5}{S5}] We-Solve research, contribution \#25, \emph{Genericity matters for odd exterior forms: a complete universal classification}; catalogue correction and stated positive generic cases, reviewed 18 October 2026. \url{https://github.com/We-Solve/research/issues/25}.
\end{description}
\label{end:1500023}
\end{document}
